%% FullCourse_10Lecs -- Lecture 9, Part 9.5a
%% Assembled by tools/lec09_assemble.py from reorg_manifest.yaml: Phase A of
%% Lec09_Reorg_Plan_2026-09-12.md as amended by Lec09_Reorg_Plan_Review_2026-09-12.md.
%% Each "% ---- <ID> · <TAG> · TIER" marker names a frame's source deck and line;
%% keep the markers until Phase B is done (lec09_assemble.py check reads them).
%% Owns: the six-step pipeline (labels verbatim; Lec10 T8 reproduces them), the 5-pillar spine (capstone scaffold)
%% Owns: restate-the-spec (verbatim sentence), five plan-grading criteria, validation checklist, V0/seed, OLG demo, tool choice, prompt specificity
%% Defers: project objects and gates at project scale -> T5b; the seven tests -> T5d

%% Shared preamble for the FullCourse_10Lecs topic decks.
%%
%% Each topic .tex \input's this file, then defines its own \title, \date
%% override (if any), \graphicspath, and \begin{document}.
%%
%% Reconstructed verbatim from the inlined copy preserved in
%% Lec10_Case_Studies/Lec10_T8_Case6_DeepRamsey.tex (lines 23-190), which is
%% explicitly annotated there as "from ../shared_preamble.tex, copied verbatim".
%% Base preamble matches MiniCourse_8hr/Slides/shared_preamble.tex; this version
%% additionally provides \sectiondivider, the conceptbox/cautionbox/examplebox/
%% takeawaybox tcolorboxes, and \compacttables used across the split decks.

\documentclass[10pt,english,aspectratio=169]{beamer}

\usepackage{ae,aecompl}
\usepackage[T1]{fontenc}
\usepackage{textcomp}
\usepackage[utf8]{inputenc}
\usepackage{babel}
\usepackage{datetime2}

\setcounter{secnumdepth}{3}
\setcounter{tocdepth}{3}

\usepackage{amsmath, amssymb, amsthm, graphicx}
\usepackage{booktabs, multirow}
\usepackage{tikz}
\usepackage{pgfplots}
\pgfplotsset{compat=1.16}
\usepackage[most]{tcolorbox}
\tcbuselibrary{skins, breakable, theorems}
\usepackage{listings}
\usepackage{xcolor}

\lstset{
  basicstyle=\ttfamily\small,
  keywordstyle=\color{blue!80!black}\bfseries,
  commentstyle=\color{green!50!black}\itshape,
  stringstyle=\color{orange!80!black},
  backgroundcolor=\color{gray!8},
  frame=single,
  rulecolor=\color{gray!40},
  breaklines=true,
  showstringspaces=false,
  numbers=none,
}

\lstdefinestyle{pythonstyle}{
    language=Python,
    basicstyle=\ttfamily\footnotesize,
    keywordstyle=\color{blue!80!black}\bfseries,
    commentstyle=\color{green!50!black}\itshape,
    stringstyle=\color{orange!80!black},
    frame=single,
    breaklines=true,
    showstringspaces=false,
    numbers=left,
    numbersep=5pt,
    numberstyle=\tiny\color{gray},
    backgroundcolor=\color{gray!8},
    rulecolor=\color{gray!40}
}

\ifx\hypersetup\undefined
  \AtBeginDocument{%
    \hypersetup{bookmarks=false,breaklinks=true,pdfborder={0 0 0},colorlinks=true,
      linkcolor=DarkText,urlcolor=RedTitle}%
  }
\else
  \hypersetup{bookmarks=false,breaklinks=true,pdfborder={0 0 0},colorlinks=true,
    linkcolor=DarkText,urlcolor=RedTitle}%
\fi

\makeatletter
\newcommand\makebeamertitle{%
  \begin{frame}[plain]
    \vspace*{1.0cm}
    \centering
    {\usebeamerfont{title}\usebeamercolor[fg]{title}\inserttitle\par}
    \vspace{1.0cm}
    {\usebeamerfont{author}\usebeamercolor[fg]{author}\insertauthor\par}
    \vspace{0.2cm}
    {\usebeamerfont{date}\usebeamercolor[fg]{date}\insertdate\par}
  \end{frame}}%
\AtBeginDocument{%
  \let\origtableofcontents=\tableofcontents
  \def\tableofcontents{\@ifnextchar[{\origtableofcontents}{\gobbletableofcontents}}
  \def\gobbletableofcontents#1{\origtableofcontents}%
}

\usetheme{metropolis}
\usefonttheme{professionalfonts}

\definecolor{LightGray}{RGB}{230,230,230}
\definecolor{RedTitle}{RGB}{0,0,200}
\definecolor{VeryLightGray}{RGB}{245,245,245}
\definecolor{DarkText}{RGB}{0,0,0}

\setbeamercolor{background canvas}{bg=white}
\setbeamercolor{title}{fg=RedTitle}
\setbeamercolor{author}{fg=DarkText}
\setbeamercolor{date}{fg=DarkText}
\setbeamercolor{frametitle}{bg=VeryLightGray,fg=DarkText}
\setbeamercolor{normal text}{fg=DarkText}
\setbeamerfont{title}{size=\large}

\setbeamersize{text margin left=5mm, text margin right=5mm}
\addtobeamertemplate{frametitle}{}{\vspace{-0.5em}}
\newcommand{\reducevspace}{\vspace{-0.5cm}}

% Shared section divider and box styles for split topic decks.
\newcommand{\sectiondivider}[2]{%
\begin{frame}[plain,noframenumbering]
\begin{center}
\vfill
{\Large\textbf{#1}\par}
\vspace{0.35cm}
{\normalsize #2\par}
\vfill
\end{center}
\end{frame}
}

\newtcolorbox{conceptbox}[1][]{
  colback=blue!3,
  colframe=RedTitle!55,
  boxrule=0.35mm,
  arc=1mm,
  left=2mm,
  right=2mm,
  top=1mm,
  bottom=1mm,
  #1
}
\newtcolorbox{cautionbox}[1][]{
  colback=red!3,
  colframe=red!50!black,
  boxrule=0.35mm,
  arc=1mm,
  left=2mm,
  right=2mm,
  top=1mm,
  bottom=1mm,
  #1
}
\newtcolorbox{examplebox}[1][]{
  colback=green!3,
  colframe=green!45!black,
  boxrule=0.35mm,
  arc=1mm,
  left=2mm,
  right=2mm,
  top=1mm,
  bottom=1mm,
  #1
}
\newtcolorbox{takeawaybox}[1][]{
  colback=VeryLightGray,
  colframe=RedTitle!55,
  boxrule=0.35mm,
  arc=1mm,
  left=2mm,
  right=2mm,
  top=1mm,
  bottom=1mm,
  #1
}
\newcommand{\compacttables}{\renewcommand{\arraystretch}{1.15}}

\setbeamertemplate{navigation symbols}{}
\setbeamertemplate{footline}{%
  \hfill%
  \usebeamercolor[fg]{page number in head/foot}%
  \usebeamerfont{page number in head/foot}%
  \insertframenumber\,/\,\inserttotalframenumber\kern1em\vskip2pt%
}

\makeatother

\author{Zhigang Feng}
% "date with time": compile timestamp so each build is identifiable.
\date{\today\ \DTMcurrenttime}


% Suppress redundant auto-generated metropolis section page; \sectiondivider frames are the dividers we keep.
\metroset{sectionpage=none}

\graphicspath{{./pic/}{./figures/}{../}}

\usetikzlibrary{positioning,arrows.meta,shapes,calc,fit,backgrounds}
\usepackage{array} % >{\raggedright\arraybackslash}p{} columns in Phase B tables

%% Lec09 shared MAP slide, v2 (September 2026 reorganization): five parts,
%% eleven decks.  v1 (the July "five acts" map) is archived as
%% _archive/five_acts_2026-07/lec09_map_v1.tex.
%%
%% Input this file AFTER ../shared_preamble.tex (needs RedTitle, VeryLightGray,
%% DarkText) and call \LecNineMap{part}{sub} inside a frame:
%%   part in {1..5} highlights that part ("you are here") and, when sub names a
%%   sub-deck letter, that deck's chip -- \LecNineMap{2}{b} is T2b, \LecNineMap{4}{}
%%   is T4;  part = 0 renders the closing reprise with every part complete (the
%%   "Your Job Description" frame in T5d).
%% Each part carries its student question and its economic twin (the delegation-
%% economics thread of Lec09_Reorg_Plan_Review_2026-09-12.md, A3).
%% Filename deliberately does not match the build glob Lec*_T*.tex, so the build
%% script never tries to compile it standalone.

% Highlight chip #1 when the requested deck key equals #2 (e.g. 2b).
\newcommand{\lecmapchip}[2]{%
  \edef\lecmaptmp{#2}%
  \ifx\lecmapkey\lecmaptmp\tikzset{#1/.style={lecchipcur}}\fi}

\newcommand{\LecNineMap}[2]{%
% TikZ option lists are comma-split before TeX conditionals run, so every
% \ifnum/\ifx is resolved here, into named styles, before the picture starts.
\edef\lecmapkey{#1#2}%
\tikzset{lecparta/.style={}, lecpartb/.style={}, lecpartc/.style={},
         lecpartd/.style={}, lecparte/.style={},
         chipTone/.style={}, chipTtwoa/.style={}, chipTtwob/.style={},
         chipTthreea/.style={}, chipTthreeb/.style={}, chipTfour/.style={},
         chipTfivea/.style={}, chipTfiveb/.style={}, chipTfivec/.style={},
         chipTfived/.style={}}%
\ifnum#1=0
  \tikzset{lecparta/.style={lecpartdone}, lecpartb/.style={lecpartdone},
           lecpartc/.style={lecpartdone}, lecpartd/.style={lecpartdone},
           lecparte/.style={lecpartdone}, lecchip/.style={lecchipbase, lecchipdone}}%
\else
  \tikzset{lecchip/.style={lecchipbase}}%
  \ifnum#1=1 \tikzset{lecparta/.style={lecpartcur}}\fi
  \ifnum#1=2 \tikzset{lecpartb/.style={lecpartcur}}\fi
  \ifnum#1=3 \tikzset{lecpartc/.style={lecpartcur}}\fi
  \ifnum#1=4 \tikzset{lecpartd/.style={lecpartcur}}\fi
  \ifnum#1=5 \tikzset{lecparte/.style={lecpartcur}}\fi
  \lecmapchip{chipTone}{1}\lecmapchip{chipTtwoa}{2a}\lecmapchip{chipTtwob}{2b}%
  \lecmapchip{chipTthreea}{3a}\lecmapchip{chipTthreeb}{3b}\lecmapchip{chipTfour}{4}%
  \lecmapchip{chipTfivea}{5a}\lecmapchip{chipTfiveb}{5b}%
  \lecmapchip{chipTfivec}{5c}\lecmapchip{chipTfived}{5d}%
\fi
\begin{center}
\begin{tikzpicture}[
    part/.style={rectangle, rounded corners, draw=black!60, fill=VeryLightGray,
        minimum width=2.62cm, minimum height=0.72cm, text width=2.5cm,
        align=center, font=\scriptsize, text=DarkText},
    lecpartcur/.style={draw=RedTitle, very thick, fill=orange!20},
    lecpartdone/.style={draw=green!45!black, thick, fill=green!10},
    lecchipbase/.style={rectangle, rounded corners=1pt, draw=black!30, fill=white,
        inner xsep=1.5pt, inner ysep=1pt, font=\tiny, text=black!70},
    lecchipcur/.style={draw=RedTitle, fill=RedTitle!12, text=RedTitle, font=\tiny\bfseries},
    lecchipdone/.style={draw=green!45!black, fill=green!8},
    q/.style={font=\tiny\itshape, text width=2.7cm, align=center, text=black!75},
    twin/.style={font=\tiny, text width=2.7cm, align=center, text=blue!40!black},
    sat/.style={rectangle, rounded corners, draw=black!45, dashed, fill=white,
        text width=5.6cm, align=center, font=\tiny, text=black!70, minimum height=0.42cm},
    barlab/.style={font=\tiny\bfseries, text=black!60},
    arr/.style={-{Stealth[length=2mm]}, thick, gray!60}
]
    % --- five part boxes ---
    \node[part, lecparta] (p1) at (0,0)     {\textbf{9.1 SEE}};
    \node[part, lecpartb] (p2) at (2.85,0)  {\textbf{9.2 UNDER THE HOOD}};
    \node[part, lecpartc] (p3) at (5.7,0)   {\textbf{9.3 BUILD}};
    \node[part, lecpartd] (p4) at (8.55,0)  {\textbf{9.4 DESIGN \& VALIDATE}};
    \node[part, lecparte] (p5) at (11.4,0)  {\textbf{9.5 RUN \& GOVERN}};
    \draw[arr] (p1) -- (p2);
    \draw[arr] (p2) -- (p3);
    \draw[arr] (p3) -- (p4);
    \draw[arr] (p4) -- (p5);
    % --- sub-deck chips ---
    \node[lecchip, chipTone]    at (0,-0.6)     {T1};
    \node[lecchip, chipTtwoa]   at (2.55,-0.6)  {T2a};
    \node[lecchip, chipTtwob]   at (3.15,-0.6)  {T2b};
    \node[lecchip, chipTthreea] at (5.4,-0.6)   {T3a};
    \node[lecchip, chipTthreeb] at (6.0,-0.6)   {T3b};
    \node[lecchip, chipTfour]   at (8.55,-0.6)  {T4};
    \node[lecchip, chipTfivea]  at (10.5,-0.6)  {T5a};
    \node[lecchip, chipTfiveb]  at (11.1,-0.6)  {T5b};
    \node[lecchip, chipTfivec]  at (11.7,-0.6)  {T5c};
    \node[lecchip, chipTfived]  at (12.3,-0.6)  {T5d};
    % --- the student question each part answers ---
    \node[q] at (0,-1.08)     {what changed,\\ and why care?};
    \node[q] at (2.85,-1.08)  {how does it run?\\ how do I start?};
    \node[q] at (5.7,-1.08)   {how do I build one,\\ and call it?};
    \node[q] at (8.55,-1.08)  {can I design one\\ and prove it works?};
    \node[q] at (11.4,-1.08)  {run a project, catch\\ mistakes, stay safe};
    % --- its economic twin ---
    \node[twin] at (0,-1.62)     {generation got cheap;\\ verification is scarce};
    \node[twin] at (2.85,-1.62)  {what am I paying for?\\ tokens, scarce context};
    \node[twin] at (5.7,-1.62)   {dividing the labour:\\ specialists, boundaries};
    \node[twin] at (8.55,-1.62)  {write the contract,\\ audit the work};
    \node[twin] at (11.4,-1.62)  {hidden action: gates,\\ monitoring, priors};
    % --- "you are here" marker above the current part ---
    \ifnum#1=1 \node[font=\scriptsize\bfseries, text=RedTitle] at (0,0.62)
        {you are here\ $\blacktriangledown$};\fi
    \ifnum#1=2 \node[font=\scriptsize\bfseries, text=RedTitle] at (2.85,0.62)
        {you are here\ $\blacktriangledown$};\fi
    \ifnum#1=3 \node[font=\scriptsize\bfseries, text=RedTitle] at (5.7,0.62)
        {you are here\ $\blacktriangledown$};\fi
    \ifnum#1=4 \node[font=\scriptsize\bfseries, text=RedTitle] at (8.55,0.62)
        {you are here\ $\blacktriangledown$};\fi
    \ifnum#1=5 \node[font=\scriptsize\bfseries, text=RedTitle] at (11.4,0.62)
        {you are here\ $\blacktriangledown$};\fi
    \ifnum#1=0 \node[font=\scriptsize\bfseries, text=green!45!black] at (5.7,0.62)
        {all five parts complete};\fi
    % --- bottom band: the three questions of the lecture ---
    \fill[blue!10]   (-1.3,-2.37) rectangle (4.15,-2.07);
    \fill[green!10]  (4.4,-2.37)  rectangle (9.85,-2.07);
    \fill[orange!15] (10.1,-2.37) rectangle (12.7,-2.07);
    \node[barlab] at (1.425,-2.22)  {HOW IT WORKS};
    \node[barlab] at (7.125,-2.22)  {HOW TO BUILD \& USE IT};
    \node[barlab] at (11.4,-2.22)   {YOUR ROLE};
    % --- dashed satellites ---
    \node[sat] at (1.9,-2.8)
        {\textbf{Appendix TA} --- the dated landscape and setup details, read on demand};
    \node[sat] at (9.9,-2.8)
        {\textbf{Lab} Parts A--B, Design Lab scenarios $\to$ \textbf{Lec10} cases (same morning)};
\end{tikzpicture}
\end{center}
}


\title[AI for Econ Research]{AI for Economic Research: Dynamic Models, Language, and Agents\\
Lecture 9.5a: Method --- The Homotopy Workflow}

\begin{document}

\makebeamertitle

% ==== Phase B · T3-F01 · TIER: CORE (route + economic twin)
\begin{frame}{Lecture 9 in One Map: The Homotopy Method}
\textbf{This deck opens the last part: how do I structure real work with the agents I just learned to build?}

\LecNineMap{5}{a}

\vspace{0.1cm}
\footnotesize\textit{Route: the homotopy idea $\to$ the six-step pipeline, step by step (with HA-Ramsey's actual opening prompts) $\to$ how to grade a plan before approving it $\to$ the runnable OLG demo $\to$ which tool for which step. The economist's question for Part 9.5: how do you manage hidden action over months --- gates, monitoring, and your own prior?}
\end{frame}


%=============================================================================
\section{The Homotopy Workflow}
%===========================================================

\sectiondivider{The Homotopy Workflow}{Build from simple to complex, validate, then extend}

% ---- T3-F02 · KEEP · TIER: READ · from T3:45
\begin{frame}{The Homotopy Idea: Solve the Easy Version First}
\textbf{In workflow design, homotopy means solving an easier version first and then extending it toward the full research task.}

\vspace{0.15cm}

\begin{columns}[T]
\begin{column}{0.55\textwidth}
\begin{itemize}
    \item Start from a version with clear intuition and known checks
    \item Use that validated baseline as the starting point for the next
    \item Attribute new bugs to the new margin you just added
    \item Never ask the agent to solve the full complicated problem in one jump
\end{itemize}
\end{column}
\begin{column}{0.42\textwidth}
\begin{center}
\begin{tikzpicture}[
    v/.style={rectangle, draw=black, fill=VeryLightGray, rounded corners,
              minimum width=1.1cm, minimum height=0.55cm,
              text width=1.0cm, align=center, font=\scriptsize},
    chk/.style={font=\tiny, text=green!50!black},
    a/.style={-{Stealth[length=1.7mm]}, thick}
]
    \node[v] (v0) {V0\\det.};
    \node[v, right=0.45cm of v0] (v1) {V1\\+risk};
    \node[v, below=0.6cm of v0] (v2) {V2\\+dist.};
    \node[v, right=0.45cm of v2] (v3) {V3\\+fiscal};
    \draw[a] (v0) -- node[above,chk]{$\checkmark$} (v1);
    \draw[a] (v1) -- ++(0,-0.4) -- node[chk]{$\checkmark$} ++(-2.0,-0.2) -- (v2);
    \draw[a] (v2) -- node[above,chk]{$\checkmark$} (v3);
\end{tikzpicture}
\end{center}
\vspace{0.05cm}
\centering\scriptsize\textit{Each step small; each step validated; if a step fails, return to the last validated version.}
\end{column}
\end{columns}

\vspace{0.1cm}

\begin{conceptbox}
\footnotesize
\textbf{Why the name:} homotopy continuation starts from a problem with a known solution and \emph{continuously deforms} it toward the target (Judd 1998, Ch.~5--7). V0 $\to$ V1 $\to$ V2 are gates along that deformation --- the same discipline, scaled up, produced HA-Ramsey (Story 1; case in Lec10 T8).
\end{conceptbox}
\end{frame}

% ---- T3-F03 · KEEP · TIER: CORE · from T3:90
% KEEP VERBATIM: Lec10 T8 reproduces these six step labels.
\begin{frame}{Homotopy in Practice: Build from Simple to Complex}
\begin{center}
\begin{tikzpicture}[
    box/.style={rectangle, draw=black, fill=VeryLightGray, rounded corners,
                minimum width=2.0cm, minimum height=0.95cm,
                text width=1.9cm, align=center, font=\small},
    arrow/.style={-{Stealth[length=2.5mm]}, thick}
]
    \node[box] (design) {1.\\Design};
    \node[box, right=0.25cm of design] (spec) {2.\\Write\\Spec};
    \node[box, right=0.25cm of spec] (algo) {3.\\Choose\\Algorithm};
    \node[box, right=0.25cm of algo] (v0) {4.\\Build\\V0};
    \node[box, right=0.25cm of v0] (val) {5.\\Validate};
    \node[box, right=0.25cm of val] (ext) {6.\\Extend};

    \draw[arrow] (design) -- (spec);
    \draw[arrow] (spec) -- (algo);
    \draw[arrow] (algo) -- (v0);
    \draw[arrow] (v0) -- (val);
    \draw[arrow] (val) -- (ext);
\end{tikzpicture}
\end{center}

\vspace{0.25cm}

\textbf{This replaces the fantasy prompt} ``solve my project.''

\begin{itemize}
    \item Start with a problem you can state and benchmark
    \item Give the agent bounded tasks with visible acceptance criteria
    \item Add one new ingredient at a time
\end{itemize}

\vspace{0.15cm}
\footnotesize\textit{The next frames walk the six steps in order --- with the HA-Ramsey case showing what each step looks like at research scale.}
\end{frame}

% ---- T3-F04 · EDIT · TIER: READ · from T3:127
\begin{frame}{Step 1 --- Design: State the Goal Like an Economist}
\textbf{Before any code, write down the task in economist language.}

\vspace{0.2cm}

\begin{itemize}
    \item Research question in one sentence
    \item Objective and constraints, or the relevant source rules
    \item Equilibrium definition, estimand, or review standard if applicable
    \item Parameters, sample, source universe, and economic motivation
    \item Outputs you will treat as evidence
\end{itemize}

\vspace{0.15cm}

\textbf{Aiyagari anchor:}
\[
\max \; \mathbb{E}_0 \sum_{t=0}^{\infty} \beta^t u(c_t)
\quad \text{s.t.} \quad
c_t + a_{t+1} = (1+r)a_t + wz_t,\; a_{t+1} \geq \underline{a}
\]

\vspace{0.1cm}

\begin{takeawaybox}
\footnotesize
\textbf{Goal template} (T2b's pen-and-paper step, made precise): \emph{``Compute X with method Y to answer Z; success looks like W.''} If you cannot fill in W, you are not ready to open the terminal.
\end{takeawaybox}
\end{frame}

% ---- T2-F23 · MOVE · TIER: CORE · from T2:747
% KEEP VERBATIM: the capstone cites these five pillar names ("your scaffold on every track", Capstone_Projects_2026.md:33).
\begin{frame}{The 5-Pillar Spine: What You Supply Before Step 2}
\textbf{The five pillars (introduced in Lec01, exercised throughout Lec06) reappear here as the agent-era research workflow.}

\vspace{0.15cm}

\begin{columns}[T]
\begin{column}{0.34\textwidth}
\footnotesize
\begin{enumerate}
    \item \textbf{Theory}
    \item \textbf{Algorithm}
    \item \textbf{Numerics}
    \item \textbf{Advanced}
    \item \textbf{Code literacy}
\end{enumerate}
\vspace{0.1cm}
\scriptsize\textit{The agent accelerates each pillar; it replaces none of them.}
\end{column}
\begin{column}{0.63\textwidth}
\footnotesize
\begin{tabular}{|p{2.1cm}|p{5.7cm}|}
\hline
\textbf{Rendering} & \textbf{One-line discipline} \\ \hline
1. Question design & Economic language; define the evidence; limiting cases first \\ \hline
2. Procedure design & The right procedure \emph{and} its acceptance criterion, before code \\ \hline
3. Data/numerical discipline & Grids, weights, crosswalks, missingness --- saved, not implicit \\ \hline
4. Domain constraints & Know when the agent's convenient default is wrong \\ \hline
5. Code literacy & Read what was built; verify it is what you meant \\ \hline
\end{tabular}
\end{column}
\end{columns}

\vspace{0.1cm}
\textbf{The throughline:} pillars 1--4 decide \emph{what should be built and what counts as success}; the machine only decides how fast. Weak pillars $=$ faster nonsense.
\end{frame}

% ---- T3-F05 · KEEP · TIER: CORE · from T3:157
% KEEP VERBATIM: Lec10 T8:540 paraphrases the restate-the-spec sentence.
\begin{frame}{Step 2 --- The Spec Is the Contract}
\textbf{The specification file is the contract between you and the agent.}

\vspace{0.2cm}

\textbf{A good specification file states four things clearly:}
\begin{enumerate}
    \item assumptions and domain rules
    \item data or source restrictions
    \item procedure and validation target
    \item required outputs and where to save them
\end{enumerate}

\vspace{0.25cm}

\textbf{Examples:} CRRA utility and convergence tolerance for Aiyagari; weights and subgroup cuts for MEPS; quote-level evidence rules for paper review; official/institutional-source rules and provenance fields for fellows data.

\vspace{0.18cm}
\textbf{Best practice:} ask the agent to restate the specification in its own words before it writes code. That catches misunderstandings when they are still cheap.
\end{frame}

% ==== Phase B · NEW:Effective Prompting: Vague vs.\ Specific · TIER: READ
\begin{frame}{Effective Prompting: Vague vs.\ Specific}
\textbf{A specific prompt names the object, the method, and the check that tells you it worked. A vague one leaves all three to the agent.}

\vspace{0.15cm}
\begin{center}
\scriptsize
\renewcommand{\arraystretch}{1.3}
\begin{tabular}{|>{\raggedright\arraybackslash}p{2.6cm}|>{\raggedright\arraybackslash}p{5.2cm}|>{\raggedright\arraybackslash}p{3.9cm}|}
\hline
\textbf{Vague} & \textbf{Specific} & \textbf{The acceptance check} \\ \hline
``Analyze this data.'' & Load \texttt{employment.csv}; compute monthly growth by sector; plot a time series with a source note & growth rates match a hand-computed sector; source and vintage printed \\ \hline
``Make it faster.'' & Vectorize the inner VFI loop with NumPy broadcasting; profile before and after & same policy function to $10^{-10}$; the profile shows the speedup \\ \hline
``Fix the regression.'' & Cluster standard errors at the level where the treatment is assigned & the clustering variable is named in the code and the table note \\ \hline
``Add some validation.'' & Add an Euler-residual check; report the max and the mean & both below $10^{-6}$ on the grid, printed every run \\ \hline
\end{tabular}
\end{center}

\vspace{0.15cm}
\begin{takeawaybox}
\footnotesize\textbf{The third column is the one agents cannot supply for you.} A specific prompt without an acceptance check is still a request for something that merely looks finished.
\end{takeawaybox}
\end{frame}

% ---- T3-F06 · KEEP · TIER: READ · from T3:178
\begin{frame}[fragile]{How HA-Ramsey Opened: Restate, Audit, Refine}
\textbf{The restate-first practice is not hypothetical --- these three prompts opened the HA-Ramsey collaboration.}

\begin{lstlisting}[basicstyle=\ttfamily\scriptsize,numbers=none]
1. "Here are the RA paper (TeX) and the RA repository.
    Restate the model and the algorithm in your own words.
    List every equation and where it is implemented."
2. "Audit the code against the document. Report every place
    where the code and the TeX disagree: scoring, bounds,
    sampling, penalties. Classify each as a correctness risk
    or a documentation gap."
3. "Propose a refined RA document + code pair in which the
    two are consistent. Do not add any new economics yet."
\end{lstlisting}

\vspace{0.05cm}

\begin{itemize}
    \item Prompt 1 tests comprehension; prompt 2 tests the \emph{seed}; prompt 3 fixes the foundation while freezing the economics
    \item ``The audit is the cheapest insurance in the whole project'' --- extending an inconsistent seed propagates the inconsistency
\end{itemize}

\vspace{0.05cm}
{\scriptsize \textit{Source: HA-Ramsey case-study package (Lec10 T8).}}
\end{frame}

% ---- T3-F07 · EDIT · TIER: READ · from T3:204
\begin{frame}{Grade the Plan Before Any Code: Five Plan-Grading Criteria}
\textbf{T2b said: start in plan mode, review, approve. Approve against \emph{what}? These five criteria.}

\vspace{0.15cm}

\begin{enumerate}
    \item \textbf{Restates your goal correctly} --- in its own words, not your words echoed back
    \item \textbf{Bounded steps} --- each step has a named deliverable small enough to review in minutes
    \item \textbf{Acceptance criterion per step} --- a benchmark, test, or check named \emph{in advance}
    \item \textbf{Stop and rollback points} --- what happens when a step fails; which validated version you return to
    \item \textbf{Smallest testable object first} --- the plan begins with V0, not with the full target
\end{enumerate}

\vspace{0.15cm}

\begin{cautionbox}
\footnotesize
\textbf{A plan without acceptance criteria is a wish.} Send it back --- a plan revision costs one message; a wrong build costs a day.
\end{cautionbox}

\vspace{0.1cm}
\footnotesize\textit{Fellows instance: one batch of 50 = a bounded step; ``every filled field carries a source URL'' = its acceptance criterion. HA-Ramsey instance: the audit prompt is criterion 3 applied to the seed itself.}
\end{frame}

% ---- T3-F08 · KEEP · TIER: READ · from T3:228
\begin{frame}[fragile]{Step 3 --- Algorithm and Pseudocode Catch Design Errors}
\textbf{Pseudocode or a phase plan is where you catch design mistakes before implementation.}

\vspace{0.2cm}

\begin{lstlisting}[style=pythonstyle,language={},basicstyle=\ttfamily\tiny,numbers=none]
Phase A: ingest or define the object
Phase B: run the chosen procedure
Phase C: save diagnostics and checks
Phase D: report only what passed validation
\end{lstlisting}

\vspace{0.15cm}

\textbf{Concrete procedures:} VFI or EGM for Aiyagari, a 3-prompt pipeline for MEPS, section decomposition for paper review, and batched metadata collection for fellows data.

\vspace{0.15cm}
\textbf{The point is not elegance.} The point is that you and the agent agree on the procedure before code details begin to hide conceptual errors.
\end{frame}

% ---- T3-F09 · KEEP · TIER: CORE · from T3:248
\begin{frame}{Step 4 --- V0: The Simplest Version With a Benchmark You Understand}
\textbf{V0 is the simplest version with a benchmark you understand --- ideally a closed form.}

\vspace{0.15cm}

\textbf{Optimal growth V0: full depreciation.} With log utility, Cobb--Douglas output $y_t = z_t k_t^{\alpha}$, and $\delta = 1$, the model has an exact closed-form solution:
\[
k_{t+1} = \alpha\beta\, z_t k_t^{\alpha}, \qquad c_t = (1-\alpha\beta)\, z_t k_t^{\alpha}
\]

\begin{itemize}
    \item Your V0 solver must reproduce this policy to numerical precision --- an \emph{absolute} benchmark, not a plausibility check
    \item Then relax $\delta < 1$: the closed form disappears, but the validated V0 has already certified the plumbing --- grids, interpolation, Euler equation, simulation
\end{itemize}

\vspace{0.15cm}

\textbf{Discipline:} do not ask the agent to add partial depreciation, risk, or equilibrium search before V0 matches the closed form.
\end{frame}

% ---- T3-F10 · KEEP · TIER: READ · from T3:268
\begin{frame}{V0 at Research Scale: What Makes a Good Seed}
\textbf{At research scale, V0 has a name: the seed. HA-Ramsey grew from a published RA solver because that seed had four properties.}

\vspace{0.2cm}

\begin{enumerate}
    \item \textbf{It works} --- reproduces its own paper's numbers, not just ``runs without error''
    \item \textbf{It is documented} --- a TeX model note in the \emph{same notation} as the code
    \item \textbf{It is modular} --- config-driven, so every later experiment is a config diff, not a code fork
    \item \textbf{It points somewhere} --- the RA model is the degenerate no-risk case of the HA target; the deformation path exists by construction
\end{enumerate}

\vspace{0.2cm}

\begin{examplebox}
\footnotesize
\textbf{A good seed is like a good prototype for manufacturing a new product:} it already works end-to-end at small scale, so scaling up is engineering and refinement --- not a fresh invention on the factory floor.
\end{examplebox}

\vspace{0.1cm}
\textbf{Same object, different scales:} the classroom V0 and the research seed play the same role. The OLG demo's \texttt{v0} (later this deck) is your training-wheels seed.
\end{frame}

% ---- T3-F11 · KEEP · TIER: READ · from T3:291
\begin{frame}{What ``V0 First'' Means Across Four Tasks}
\begin{center}
\small
\begin{tabular}{|p{2.2cm}|p{3.3cm}|p{5.2cm}|}
\hline
\textbf{Task} & \textbf{V0} & \textbf{First validation target} \\ \hline
Aiyagari model & deterministic household problem & recover simple savings logic and boundary behavior \\ \hline
MEPS pipeline & 3-year pilot panel & weighted rates line up with an external benchmark \\ \hline
Paper review skill & one section or theorem block & quotes are real and comments are specific \\ \hline
Fellows dataset & seed roster only & source rules and provenance fields are actually recorded \\ \hline
\end{tabular}
\end{center}

\vspace{0.25cm}
\textbf{Why economists should care:} V0 keeps the first error interpretable. You want the first failure to be obvious, not hidden inside a giant workflow.
\end{frame}

% ---- T3-F12 · EDIT · TIER: CORE · from T3:308
\begin{frame}{Steps 5--6 --- Validate, Then Extend One Margin at a Time}
\begin{center}
\small
\begin{tabular}{|c|p{3.5cm}|p{4.0cm}|}
\hline
\textbf{Version} & \textbf{New ingredient} & \textbf{Required benchmark} \\ \hline
V0 & Deterministic household & analytical cases \\ \hline
V1 & One new model, data, or review margin & setting the new margin to zero recovers V0 \\ \hline
V2 & Saved diagnostics and state tracking & mass sums, benchmarks, quote checks, or source logs behave sensibly \\ \hline
V3 & Full target workflow & final artifact survives validation and interpretation \\ \hline
\end{tabular}
\end{center}

\vspace{0.2cm}

\textbf{Why this works with agents:} when a new version breaks, both you and the agent know where to look.

\vspace{0.1cm}
\footnotesize\textit{Hold the gate idea: in T5b it returns as a management principle --- a quality gate on every intermediate good, not just on code versions.}
\end{frame}

% ---- T2-F25 · MOVE · TIER: READ · from T2:802
\begin{frame}{Validation Is Your Referee Checklist}
\textbf{Every iteration needs both workflow checks and economic checks.}

\vspace{0.25cm}

\begin{columns}[T]
\begin{column}{0.48\textwidth}
\textbf{Workflow checks}
\begin{itemize}
    \item Convergence achieved
    \item No NaN or Inf values
    \item Mappings, weights, or source logs saved where expected
    \item Diagnostics, residuals, or quote checks pass
    \item Script reproduces the same output twice
\end{itemize}
\end{column}
\begin{column}{0.48\textwidth}
\textbf{Economic checks}
\begin{itemize}
    \item Limiting cases recover known results
    \item Comparative statics, subgroup patterns, or review claims make sense
    \item Magnitudes and source coverage are plausible
    \item Inference uses the right standard errors, sample, and restrictions
    \item You can explain the result in plain economic language
\end{itemize}
\end{column}
\end{columns}

\vspace{0.25cm}

\textbf{If you cannot defend the output in a seminar, you are not done.}
\end{frame}

% ---- T3-F13 · KEEP · TIER: READ · from T3:329
\begin{frame}{Three Layers Grow Together: Model, Algorithm, Code}
\textbf{A research project is not one artifact but three, and the homotopy steps move all three in lockstep.}

\vspace{0.15cm}

\begin{center}
\begin{tikzpicture}[
    layer/.style={rectangle, rounded corners, draw=RedTitle!70, fill=blue!5,
        minimum width=3.1cm, minimum height=0.95cm, text width=3.0cm,
        align=center, font=\small},
    arr/.style={{Stealth[length=2mm]}-{Stealth[length=2mm]}, thick, gray!70},
    back/.style={-{Stealth[length=2mm]}, thick, dashed, green!45!black}
]
    \node[layer] (model) at (0,0) {\textbf{Model}\\ \scriptsize TeX: economics, timing, constraints};
    \node[layer] (algo) at (4.5,0) {\textbf{Algorithm}\\ \scriptsize pseudocode: solution method, penalties};
    \node[layer] (code) at (9.0,0) {\textbf{Code}\\ \scriptsize PyTorch/Python: what actually runs};
    \draw[arr] (model) -- (algo);
    \draw[arr] (algo) -- (code);
    \draw[back] (code.south) to[bend right=20] node[below, font=\tiny, text=green!45!black]
        {step back anytime: check, validate, revert} (model.south);
\end{tikzpicture}
\end{center}

\vspace{0.1cm}

\begin{itemize}
    \item \textbf{Forward:} you grow from the current version toward the final version gradually --- a new model element forces a new algorithm forces new code, one margin at a time
    \item \textbf{Backward:} every earlier version is preserved and validated, so you can always step back to a simpler one --- to check and validate a result, to isolate where something broke, or to revert to the last working version
    \item \textbf{The agent's job here:} keep TeX, pseudocode, and code describing \emph{the same economy} at every version --- drift between layers is where silent errors live
\end{itemize}

\vspace{0.05cm}
{\scriptsize \textit{Source: HA-Ramsey case-study package (Lec10 T8).}}
\end{frame}

% ---- T3-F14 · KEEP · TIER: READ · from T3:364
\begin{frame}{Why This Beats One-Shot Generation}
\textbf{Four reasons the homotopy workflow is especially effective with agents:}

\begin{enumerate}
    \item \textbf{Agents excel at well-specified, bounded tasks.} \\
    ``Solve my model'' fails. ``Implement V0 per this pseudocode and validate against these 3 benchmarks'' succeeds.
    \medskip
    \item \textbf{Starting simple gives both you and the agent a working codebase.} \\
    V1 inherits validated V0 code. The agent builds on tested functions, not a blank file.
    \medskip
    \item \textbf{Bug attribution is immediate.} \\
    If V2 breaks, the bug is in V2's new lines, not the validated inherited ones.
    \medskip
    \item \textbf{Context builds naturally.} \\
    By V3, the agent has seen the spec, pseudocode, three validated versions, and every correction. \emph{Deep context} --- not a 10-page prompt.
\end{enumerate}

\textbf{Discipline:} resist the urge to skip ahead. A working V0 takes 30 minutes; debugging a broken V2 because V0 was never validated takes a day. This is also how good RAs are trained: simple task, verify, widen scope.
\end{frame}


%=============================================================================
\section{OLG Five-Step Demo}
%===========================================================

\sectiondivider{OLG Five-Step Demo}{A concrete homotopy workflow for computational macro}

% ---- T3-F15 · KEEP · TIER: LAB · from T3:393
\begin{frame}{Demo: Folder Structure Is the Method}
\small
\begin{tabular}{p{4.0cm}p{7.7cm}}
\toprule
\textbf{Folder or file} & \textbf{Teaching role} \\
\midrule
\texttt{stages/} & five durable stages: model, equilibrium, algorithm, pseudocode, implementation/validation \\
\texttt{prompts/v*\_to\_v*.md} & one bounded terminal-agent session per version transition \\
\texttt{versions/v0--v5} & validated working code preserved at every step \\
\texttt{tests/} & one smoke test per version; failures identify the new margin \\
\texttt{build/*.json} & reproducibility reports from terminal runs \\
\bottomrule
\end{tabular}

\vspace{0.25cm}
\textbf{Demo path:} \texttt{demos/M4\_olg\_5step\_demo/} (local copy in this lecture; mirror of MiniCourse\_8hr \texttt{M4\_olg\_5step\_demo} minus \texttt{.venv}/\texttt{build/}).

\vspace{0.1cm}
\footnotesize\textit{Note the pattern: the HA-Ramsey \texttt{Case\_Study/} package (Lec10 T8) is this same folder discipline at paper scale --- seed, versions, notes, audits, all preserved.}
\end{frame}

% ---- T3-F16 · KEEP · TIER: LAB · from T3:414
\begin{frame}{Version Ladder: One Economic Margin at a Time}
\scriptsize
\begin{tabular}{p{0.8cm}p{5.0cm}p{5.8cm}}
\toprule
\textbf{Ver.} & \textbf{Adds} & \textbf{Validation gate} \\
\midrule
V0 & 3 cohorts, 1 asset, 2-state TFP, one NN policy & deterministic steady-state collapse; procyclical capital; RMS Euler residual below threshold \\
V1 & 7 cohorts and lifecycle labor & V0 logic preserved; middle-aged consumption peak appears \\
V2 & 4-state Markov TFP & collapse back to two states reproduces V1 \\
V3 & bonds, price, clearing, complementarity & zero bond-loss weight reproduces V2; positive late-life holdings \\
V4 & capital adjustment cost & zero adjustment cost reproduces V3; capital path smooths \\
V5 & stabilizing homotopy phases & residuals fall by phase; final Euler error target met \\
\bottomrule
\end{tabular}

\vspace{0.25cm}
\textbf{Research lesson:} the agent never receives ``solve the big model.'' It receives one validated deformation at a time.
\end{frame}

% ==== Phase B · T3-F17 · TIER: LAB (README launch order)
\begin{frame}[fragile]{Demo Launch and Classroom Contract}
\begin{lstlisting}[basicstyle=\ttfamily\scriptsize, numbers=none]
cd source/Lec09_Agentic_AI/demos/M4_olg_5step_demo
pip install -e .
jupyter notebook demo.ipynb          # Run All: V0 -> V5 with each gate

python3 -m unittest discover -s tests  # separate smoke check
\end{lstlisting}

\vspace{0.2cm}
\begin{itemize}
  \item \textbf{What the demo proves:} homotopy preserves a working baseline while complexity rises.
  \item \textbf{What to inspect:} prompts, version folders, test names, validation messages, and regenerated reports.
  \item \textbf{Failure recovery:} if V3 breaks, return to validated V2 and add only the bond-market-clearing layer.
\end{itemize}
\end{frame}


%=============================================================================
\section{Tool Choice}
%===========================================================

\sectiondivider{Tool Choice}{Match the tool to the pipeline step}

% ==== Phase B · T3-F18 · TIER: READ (harness-neutral)
\begin{frame}{Tool Choice per Pipeline Step}
\textbf{The six-step pipeline also tells you which tool to open.}

\vspace{0.15cm}

\begin{center}
\small
\begin{tabular}{|p{2.4cm}|p{3.4cm}|p{5.2cm}|}
\hline
\textbf{Tool type} & \textbf{Pipeline steps} & \textbf{Economic example} \\ \hline
Web chat & 1--2: design, draft the spec & clarify identification logic, draft a model summary \\ \hline
IDE assistant & single-file edits inside a step & vectorize one function, explain one traceback \\ \hline
Terminal agent (Claude Code, Codex CLI, Gemini CLI) & 4--6: build, validate, extend & restructure the project, run scripts, generate diagnostics, iterate \\ \hline
\end{tabular}
\end{center}

\vspace{0.2cm}

\textbf{Why the terminal agent becomes the workhorse:} it runs the real code on the real project, while leaving you in the approval and validation loop --- enough autonomy to save time, not so much that you stop checking the economics.

\vspace{0.1cm}
{\footnotesize The OLG demo ships both \texttt{CLAUDE.md} and \texttt{AGENTS.md}: one project, briefed for more than one harness.}
\end{frame}

% ==== Phase B · T3-F19 · TIER: READ (bridge)
\begin{frame}{Bridge: One Ladder vs.\ the Whole Factory}
\textbf{T5a disciplined the growth of one artifact. A real project is a factory: many artifacts, many sessions, and a team of the agents you built in Part 9.3.}

\vspace{0.25cm}

\begin{itemize}
    \item You now have the \emph{method}: design, spec, algorithm, V0, validate, extend --- with a plan graded before any code
    \item But a paper is model + algorithm + code + data + reviews + write-up, worked across weeks of sessions
    \item The next question: \emph{``how do I run the whole project?''} --- T5b's answer is a role change: the researcher as \textbf{project manager}, and the project files (instructions, skills, agents, rules) as the manager's instruments
\end{itemize}

\vspace{0.3cm}
\footnotesize\textit{Arc: 9.1 See $\to$ 9.2 Under the Hood $\to$ 9.3 Build $\to$ 9.4 Design \& Validate $\to$ \textbf{9.5 Run \& Govern}.}
\end{frame}

% ==== Phase B · NEW:Appendix: A Concrete Aiyagari Spec and Pseudocode · TIER: READ
\begin{frame}[fragile]{Appendix: A Concrete Aiyagari Spec and Pseudocode}
\textbf{Steps 1--2 for a model this course already uses: the spec the agent reads every session, and pseudocode agreed before any Python.}

\vspace{0.1cm}
\begin{columns}[T]
\begin{column}{0.48\textwidth}
\footnotesize\textbf{\texttt{model\_spec.md} --- a minimal version}
\begin{enumerate}\scriptsize
    \item \textbf{Household:} CRRA utility, $\beta = 0.96$, $\sigma = 2$
    \item \textbf{Income:} AR(1) in logs, $\rho = 0.9$, $\sigma_\varepsilon = 0.2$, Tauchen with 7 states
    \item \textbf{Asset grid:} 200 points, spaced densely near the borrowing limit
    \item \textbf{Method:} VFI with Howard improvement
    \item \textbf{Convergence:} sup-norm on $V$ below $10^{-8}$
    \item \textbf{Equilibrium:} $r$ clears the asset market, by bisection
    \item \textbf{Output:} $c(a,z)$, $a'(a,z)$, the stationary distribution, the Gini coefficient
\end{enumerate}
{\scriptsize Then ask the agent to restate the model in its own words.}
\end{column}
\begin{column}{0.5\textwidth}
\footnotesize\textbf{Pseudocode for V0 (deterministic income)}
\begin{lstlisting}[style=pythonstyle,language={},basicstyle=\ttfamily\tiny,numbers=none]
grid a[1..N]; guess V_old = u(r*a + w) / (1 - beta)
repeat until ||V_new - V_old|| < tol:
    for each a_i:
        for each a' with c = (1+r)*a_i + w - a' > 0:
            value = u(c) + beta * interp(V_old, a')
        V_new[i]    = max over a'
        a_policy[i] = argmax over a'
    every 25 iterations: 50 Howard steps
return V_new, a_policy, c = (1+r)*a + w - a_policy
\end{lstlisting}
{\scriptsize Note the \texttt{c > 0} guard and the stopping rule --- the two things the lab's flawed solver leaves out (T3b, T4).}
\end{column}
\end{columns}
\end{frame}

% ==== Phase B · NEW:Appendix: Three Analytical V0 Benchmarks · TIER: READ
\begin{frame}{Appendix: Three Analytical V0 Benchmarks}
\textbf{Before V1, the deterministic-income V0 must reproduce three cases you can derive on paper. If it fails one, nothing built on it can be trusted.}

\vspace{0.2cm}
\begin{center}
\small
\renewcommand{\arraystretch}{1.35}
\begin{tabular}{|>{\raggedright\arraybackslash}p{2.6cm}|>{\raggedright\arraybackslash}p{4.4cm}|>{\raggedright\arraybackslash}p{4.4cm}|}
\hline
\textbf{Calibration} & \textbf{What theory says} & \textbf{What V0 must show} \\ \hline
$\beta(1+r) = 1$ & the Euler equation holds with constant consumption & $a'(a) = a$: the policy lies on the 45-degree line \\ \hline
$\beta(1+r) < 1$ & impatience outweighs the return: dissave & assets run down to the borrowing limit \\ \hline
$\beta(1+r) > 1$ & the return outweighs impatience: save & assets keep growing toward the top of the grid \\ \hline
\end{tabular}
\end{center}

\vspace{0.2cm}
\begin{takeawaybox}
\footnotesize\textbf{If V0 fails:} the Bellman equation, the grid, or the stopping rule is wrong --- and every later version inherits it. A working V0 takes an afternoon; debugging a V2 built on a broken V0 takes a week.
\end{takeawaybox}
\end{frame}

%===========================================================
\end{document}
